%----------------------------------------------------------------------------------------
%	VARIOUS REQUIRED PACKAGES
%----------------------------------------------------------------------------------------

\usepackage{tikz} % Required for drawing custom shapes
%\usepackage{booktabs} % Required for nicer horizontal rules in tables
\usepackage{eso-pic} % Required for specifying an image background in the title page


%----------------------------------------------------------------------------------------
%	SECTIONS AND SUBSECTIONS
%----------------------------------------------------------------------------------------
%\titleformat*{\section}{\Huge}
%\titleformat*{\subsection}{\Large}


% Chapter and section formatting 
\usepackage[rigidchapters,explicit]{titlesec}
\setlength{\fboxsep}{5pt}
\newlength{\titulolength}
\setlength{\titulolength}{\textwidth} 
\addtolength{\titulolength}{-2\fboxsep}
\addtolength{\titulolength}{-2\fboxrule}
\titleformat{\chapter}[display]
{\bfseries\sffamily}
{\filleft \mdseries \slshape \LARGE \color{blueceu} Práctica \Huge\thechapter}
{0.5ex}
{\colorbox{blueceu}{\parbox[t]{\titulolength}{\rule{0mm}{10mm}\Huge\sffamily\filleft\textcolor{white}{#1}}}}
\titleformat{\section}{\normalfont\Large\bfseries\color{blueceu}}{\arabic{section}}{1em}{#1}
\titleformat{\subsection}{\normalfont\large\bfseries\color{blueceu}}{\arabic{section}.\arabic{subsection}}{1em}{#1}
\titleformat{\subsubsection}{\normalfont\normalsize\bfseries\color{blueceu}}{\arabic{section}.\arabic{subsection}.\arabic{subsubsection}}{0.5em}{#1}
\titlespacing{\chapter}{0pt}{0pt}{30ex}
\titlespacing{\section}{0pt}{8ex}{4ex}

%----------------------------------------------------------------------------------------
%	MAIN TABLE OF CONTENTS
%----------------------------------------------------------------------------------------

% \usepackage{titletoc} % Required for manipulating the table of contents
% 
% \contentsmargin{0cm} % Removes the default margin
% % Chapter text styling
% \titlecontents{chapter}[1.25cm] % Indentation
% {\addvspace{15pt}\large\bfseries} % Spacing and font options for chapters
% {\color{gray}\contentslabel[\Large\thecontentslabel.]{1.25cm}} % Chapter number
% {}  
% {\color{gray}\bfseries\hfill\thecontentspage} % Page number
% % Section text styling
% \titlecontents{section}[1.25cm] % Indentation
% {\addvspace{5pt}\bfseries} % Spacing and font options for sections
% {\contentslabel[\thecontentslabel]{1.25cm}} % Section number
% {}
% {\hfill\color{black}\thecontentspage} % Page number
% []
% % Subsection text styling
% \titlecontents{subsection}[1.25cm] % Indentation
% {\addvspace{1pt}\small} % Spacing and font options for subsections
% {\contentslabel[\thecontentslabel]{1.25cm}} % Subsection number
% {}
% {\;\titlerule*[.5pc]{.}\;\thecontentspage} % Page number
% [] 

%----------------------------------------------------------------------------------------
%	MINI TABLE OF CONTENTS IN CHAPTER HEADS
%----------------------------------------------------------------------------------------

% % Section text styling
% \titlecontents{lsection}[0em] % Indendating
% {\footnotesize\sffamily} % Font settings
% {}
% {}
% {}
% 
% % Subsection text styling
% \titlecontents{lsubsection}[.5em] % Indentation
% {\normalfont\footnotesize\sffamily} % Font settings
% {}
% {}
% {}
 
%----------------------------------------------------------------------------------------
%	PAGE HEADERS
%----------------------------------------------------------------------------------------

% Headings and footers
\usepackage{fancyhdr}
\pagestyle{fancy}
\fancyhead{}
\renewcommand{\chaptermark}[1]{\markboth{\thechapter.\ #1}{}}
\fancyhead[LE]{\scshape Bioestadística Aplicada con R y RKTeaching}
\fancyhead[RO]{\sffamily\slshape \leftmark}
\renewcommand{\headrulewidth}{0pt}
\renewcommand{\floatpagefraction}{.8}
\renewcommand{\textfraction}{.1} 


% \usepackage{fancyhdr} % Required for header and footer configuration
% 
% \pagestyle{fancy}
% \renewcommand{\chaptermark}[1]{\markboth{\color{grayceu}\sffamily\normalsize\textit{#1}}{}} % Chapter text font
% % settings
% %\renewcommand{\sectionmark}[1]{\markright{\color{grayceu}\sffamily\normalsize\textit{\thesection\hspace{5pt}#1}}{}} %
% % Section text font settings
% \fancyhf{} \fancyfoot[LE,RO]{\normalsize\thepage} % Font setting for the page number in the footer
% \fancyhead[RE]{{\color{grayceu}\sffamily\normalsize\textit{Pr�cticas de Bioestad�stica con R y RKTeaching}}} % Print the nearest
% % section name on the left side of odd pages
% \fancyhead[LO]{\leftmark} % Print the current chapter name on the right side of even pages
% \renewcommand{\headrulewidth}{0pt} % Removes the rule in the header
% \addtolength{\headheight}{2.5pt} % Increase the spacing around the header slightly
% \renewcommand{\footrulewidth}{0pt} % Removes the rule in the footer
% \fancypagestyle{plain}{\fancyhead{}\renewcommand{\headrulewidth}{0pt}} % Style for when a plain pagestyle is specified
% 
% % Removes the header from odd empty pages at the end of chapters
% \makeatletter
% \renewcommand{\cleardoublepage}{
% \clearpage\ifodd\c@page\else
% \hbox{}
% \vspace*{\fill}
% \thispagestyle{empty}
% \newpage
% \fi}


%----------------------------------------------------------------------------------------
%	THEOREM STYLES
%----------------------------------------------------------------------------------------

\usepackage{amsthm} % For including math equations, theorems, symbols, etc

\newcommand{\intoo}[2]{\mathopen{]}#1\,;#2\mathclose{[}}
\newcommand{\ud}{\mathop{\mathrm{{}d}}\mathopen{}}
\newcommand{\intff}[2]{\mathopen{[}#1\,;#2\mathclose{]}}
\newtheorem{notation}{Notation}[chapter]

\newtheoremstyle{blueceu} % Theorem style name
{7pt} % Space above
{7pt} % Space below
{\normalfont} % Body font
{} % Indent amount
{\small\bf\sffamily\color{blueceu}} % Theorem head font
{\;\;} % Punctuation after theorem head
{0.25em} % Space after theorem head
{\small\sffamily\color{blueceu}\thmname{#1}\thmnumber{\@ifnotempty{#1}{ }\@upn{#2}} % Theorem text (e.g. Theorem 2.1)
\thmnote{\ {\the\thm@notefont\sffamily\bfseries\color{blueceu}--- #3.}}} % Optional theorem note
\renewcommand{\qedsymbol}{$\blacksquare$} % Optional qed square

\newtheoremstyle{blacknumex} % Theorem style name
{7pt} % Space above
{7pt} % Space below
{\normalfont} % Body font
{} % Indent amount
{\small\bf\sffamily} % Theorem head font
{\;\;} % Punctuation after theorem head
{0.25em} % Space after theorem head
{\small\sffamily{\tiny\ensuremath{\blacksquare}}\ \thmname{#1}\thmnumber{\@ifnotempty{#1}{ }\@upn{#2}} % Theorem text (e.g. Theorem 2.1)
\thmnote{\ {\the\thm@notefont\sffamily\bfseries--- #3.}}} % Optional theorem note

\newtheoremstyle{blacknum} % Theorem style name
{7pt} % Space above
{7pt} % Space below
{\normalfont} % Body font
{} % Indent amount
{\small\bf\sffamily} % Theorem head font
{\;\;} % Punctuation after theorem head
{0.25em} % Space after theorem head
{\small\sffamily\thmname{#1}\thmnumber{\@ifnotempty{#1}{ }\@upn{#2}} % Theorem text (e.g. Theorem 2.1)
\thmnote{\ {\the\thm@notefont\sffamily\bfseries--- #3.}}} % Optional theorem note

\newtheoremstyle{example} % Theorem style name
{7pt} % Space above
{7pt} % Space below
{\normalfont} % Body font
{} % Indent amount
{\small\bf\sffamily} % Theorem head font
{\;\;} % Punctuation after theorem head
{0.25em} % Space after theorem head
{\small\sffamily\thmname{#1} % Theorem text (e.g. Theorem 2.1)
\thmnote{\ {\the\thm@notefont\sffamily\bfseries--- #3.}}} % Optional theorem note

\newtheoremstyle{indication} % Theorem style name
{-5pt} % Space above
{7pt} % Space below
{\normalfont} % Body font
{-28pt} % Indent amount
{\bf} % Theorem head font
{\kern-11.5pt} % Punctuation after theorem head
{20pt} % Space after theorem head
{\begin{tikzpicture}
\draw[anchor=west] (-26pt,10pt) node [fill=blueceu,opacity=1,inner
sep=1pt]{\Large\bfseries\textcolor{white}{\vphantom{plPQq}\makebox[17pt][c]{\includegraphics[scale=0.3, trim=0 5mm 0 0]{img/bombilla}}}};
\end{tikzpicture}} % Theorem text (e.g. Theorem 2.1)

\makeatother

% Defines the theorem text style for each type of theorem to one of the three styles above
\theoremstyle{ocrenum}
\newtheorem{proposicion}{Proposición}[chapter]
\newtheorem{problema}{Problema}[chapter]
\newtheorem{ejercicio}{Ejercicio}[chapter]
\theoremstyle{blacknumex}
\newtheorem{ejemplo}{Ejemplo}[chapter]
%\theoremstyle{blacknum}
\theoremstyle{blueceu}
\newtheorem{vocabulario}{Vocabulario}[chapter]
\newtheorem{definicion}{Definición}[chapter]
\newtheorem{teorema}{Teorema}[chapter]
\newtheorem{corolario}{Corolario}[chapter]
\theoremstyle{indication}
\newtheorem{indicacionT}{Indicación}
%\theoremstyle{example}
%\newtheorem{ejemplo}{Ejemplo}

%----------------------------------------------------------------------------------------
%	DEFINITION OF COLORED BOXES
%----------------------------------------------------------------------------------------

\RequirePackage[framemethod=default]{mdframed} % Required for creating the theorem, definition, exercise and corollary boxes

% Theorem box
\newmdenv[skipabove=7pt,
skipbelow=7pt,
backgroundcolor=black!5,
linecolor=ocre,
innerleftmargin=5pt,
innerrightmargin=5pt,
innertopmargin=5pt,
leftmargin=0cm,
rightmargin=0cm,
innerbottommargin=5pt]{tBox}

% Exercise box	  
\newmdenv[skipabove=7pt,
skipbelow=7pt,
rightline=false,
leftline=true,
topline=false,
bottomline=false,
backgroundcolor=ocre!10,
linecolor=ocre,
innerleftmargin=5pt,
innerrightmargin=5pt,
innertopmargin=5pt,
innerbottommargin=5pt,
leftmargin=0cm,
rightmargin=0cm,
linewidth=4pt]{eBox}	

% Definition box
\newmdenv[skipabove=10pt,
skipbelow=10pt,
rightline=false,
leftline=true,
topline=false,
bottomline=false,
linecolor=ocre,
innerleftmargin=5pt,
innerrightmargin=5pt,
innertopmargin=0pt,
leftmargin=0cm,
rightmargin=0cm,
linewidth=4pt,
innerbottommargin=0pt]{dBox}	

% Corollary box
\newmdenv[skipabove=7pt,
skipbelow=7pt,
rightline=false,
leftline=true,
topline=false,
bottomline=false,
linecolor=gray,
backgroundcolor=black!5,
innerleftmargin=5pt,
innerrightmargin=5pt,
innertopmargin=5pt,
leftmargin=0cm,
rightmargin=0cm,
linewidth=4pt,
innerbottommargin=5pt]{cBox}	

% Indication box
\newmdenv[skipabove=7pt,
skipbelow=7pt,
rightline=false,
leftline=true,
topline=false,
bottomline=false,
linecolor=blueceu,
backgroundcolor=black!5,
innerleftmargin=5pt,
innerrightmargin=5pt,
innertopmargin=5pt,
leftmargin=0pt,
rightmargin=0pt,
linewidth=4pt,
innerbottommargin=5pt]{iBox}		  
		  			  


% Creates an environment for each type of theorem and assigns it a theorem text style from the "Theorem Styles" section above and a colored box from above
\newenvironment{theorem}{\begin{tBox}\begin{theoremeT}}{\end{theoremeT}\end{tBox}}
\newenvironment{exercise}{\begin{eBox}\begin{exerciseT}}{\hfill{\color{ocre}\tiny\ensuremath{\blacksquare}}\end{exerciseT}\end{eBox}}				  
\newenvironment{definition}{\begin{dBox}\begin{definitionT}}{\end{definitionT}\end{dBox}}	
\newenvironment{example}{\begin{exampleT}}{\hfill{\tiny\ensuremath{\blacksquare}}\end{exampleT}}		
\newenvironment{corollary}{\begin{cBox}\begin{corollaryT}}{\end{corollaryT}\end{cBox}}	
\newenvironment{indicacion}{\begin{iBox}\begin{indicacionT}}{\end{indicacionT}\end{iBox}}	

%----------------------------------------------------------------------------------------
%	REMARK ENVIRONMENT
%----------------------------------------------------------------------------------------

%\newenvironment{remark}{\par\vskip10pt\small % Vertical white space above the remark and smaller font size
%\begin{list}{}{
%\leftmargin=35pt % Indentation on the left
%\rightmargin=25pt}\item\ignorespaces % Indentation on the right
%\makebox[-2.5pt]{\begin{tikzpicture}[overlay]
%\node[draw=ocre!60,line width=1pt,circle,fill=ocre!25,font=\sffamily\bfseries,inner sep=2pt,outer sep=0pt] at (-15pt,0pt){\textcolor{ocre}{R}};\end{tikzpicture}} % Orange R in a circle
%\advance\baselineskip -1pt}{\end{list}\vskip5pt} % Tighter line spacing and white space after remark

% %----------------------------------------------------------------------------------------
% %	SECTION NUMBERING IN THE MARGIN
% %----------------------------------------------------------------------------------------
% 
% \makeatletter
% %\renewcommand{\@seccntformat}[1]{\llap{\textcolor{grayceu}{\csname the#1\endcsname.}\hspace{1em}}}                    
% \renewcommand{\section}{
% 	\@startsection{section}{1}{\z@}
% 	{-4ex \@plus -1ex \@minus -.4ex}{1ex \@plus.2ex }
% 	{\normalfont\LARGE}
% }
% \renewcommand{\subsection}{
% 	\@startsection{subsection}{2}{\z@}
% 	{-3ex \@plus -0.1ex \@minus -.4ex}{0.5ex \@plus.2ex }
% 	{\normalfont\Large}
% }
% \renewcommand{\subsubsection}{
% 	\@startsection{subsubsection}{3}{\z@}
% 	{-2ex \@plus -0.1ex \@minus -.2ex}{0.2ex \@plus.2ex }
% 	{\normalfont\large}
% }                        
% \renewcommand{\paragraph}{
% 	\@startsection{paragraph}{4}{\z@}
% 	{-2ex \@plus-.2ex \@minus .2ex}{0.1ex}
% 	{\normalfont\bfseries}
% }
% 
% \def\thechapter{\arabic{chapter}}
% \def\thesection{\thechapter.\arabic{section}.}
% \def\thesubsection{\thesection\arabic{subsection}.}
% \def\thesubsubsection{\thesubsection\arabic{section}.}
% 
% %----------------------------------------------------------------------------------------
% %	CHAPTER HEADINGS
% %----------------------------------------------------------------------------------------
% %\newcommand{\thechapterimage}{}
% %\newcommand{\chapterimage}[1]{\renewcommand{\thechapterimage}{#1}}
% \def\@makechapterhead#1{
% \thispagestyle{empty}
% {
% \ifnum \c@secnumdepth >\m@ne
% \if@mainmatter
% \startcontents
% \begin{tikzpicture}[remember picture,overlay]
% % \node at (current page.north west)
% % {\begin{tikzpicture}[remember picture,overlay]
% % 
% % %\node[anchor=north west] at (-4pt,-27.5mm) {\includegraphics[width=\paperwidth]{\thechapterimage}};
% % %Commenting the 3 lines below removes the small contents box in the chapter heading
% % % \draw[fill=white,opacity=.6] (30mm,-30mm) rectangle (9cm,-9cm);
% % % \node[anchor=north west] at (30mm,-30mm) {\parbox[t][9cm][t]{5.5cm}{\huge\bfseries\flushleft \printcontents{l}{1}{\setcounter{tocdepth}{2}}}};
% % % \draw[anchor=west] (5cm,-11cm) node [rounded corners=25pt,fill=white,opacity=.7,inner sep=15.5pt]{\huge\sffamily\bfseries\textcolor{black}{\vphantom{plPQq}\makebox[20cm]{}}};
% \draw[anchor=north west, fill=black] (0mm,50mm) rectangle (1pt,10mm);
% \draw[anchor=north west] (0cm,20mm) node {\scalebox{1.5}{\Huge{\textcolor{grayceu}\thechapter}}}; 
% %\draw[anchor=west] (0cm,0cm) node {\Huge{#1}}; 
% \end{tikzpicture}
% \par\vspace*{2cm}
% {\Huge #1}
% \par\vspace*{1cm}
% % \end{tikzpicture}}\par\vspace*{230\p@}
% }}
% \def\@makeschapterhead#1{
% \thispagestyle{empty}
% {
% \ifnum \c@secnumdepth >\m@ne
% \if@mainmatter
% \startcontents
% \begin{tikzpicture}[remember picture,overlay]
% % \node at (current page.north west)
% % {\begin{tikzpicture}[remember picture,overlay]
% % \node[anchor=north west] at (-4pt,-27.5mm) {\includegraphics[width=\paperwidth]{\thechapterimage}};
% % \draw[anchor=west] (5cm,-11cm) node [rounded corners=25pt,fill=white,opacity=.7,inner sep=15.5pt]{\huge\sffamily\bfseries\textcolor{black}{\vphantom{plPQq}\makebox[20cm]{}}};
% % \draw[anchor=west] (5cm,-11cm) node [rounded corners=25pt,inner sep=15.5pt]{\huge\sffamily\bfseries\textcolor{black}{#1\vphantom{plPQq}\makebox[20cm]{}}};
% % \end{tikzpicture}};
% % \end{tikzpicture}}\par\vspace*{230\p@}
% %\draw[anchor=west] (0cm,0cm) node {\Huge{#1}}; 
% \end{tikzpicture}
% {\Huge #1}
% \par\vspace*{1cm}
% }}
% \makeatother

